Our aim is to solve the Hamiltonian equations of motion \ref{eq: HMC_DGLs}, which are at the heart of the HMC algorithm, numerically. This chapter 
describes different methods used for this task. 
\section{Leapfrog Integration}
In Section \ref{sec: HMC}, we have discussed that the detailed balance condition for the HMC algorithm imposes the following two 
requirements on the molecular dynamics trajectory: 
\begin{itemize}
    \item Invariance of the integration measure under Hamiltonian evolution in the $(\phi, \pi)$ phase space: 
    \begin{equation}
        [\mathcal{D}\phi] [\mathcal{D}\pi] = [\mathcal{D}\phi(\tau)] [\mathcal{D}\pi(\tau)]
    \end{equation}
    \item Reversibility of the trajectory, evolving forward in fictitious time $\tau$ and then reversing the momenta exactly retraces the trajectory.
\end{itemize}
The leapfrog integration method fulfills these requirements, as will be shown in the following. Leapfrog is a second-order method, compared to the very-well known first-order Euler method. It provides a higher accuracy while maintaining measure invariance and reversibility.  
\\
In the leapfrog integration method, the field $\phi$ and the conjugate momenta $\pi$ are evolved in $n$ steps of length $\varepsilon$. The conjugate momenta start with a half-step of $\frac{\varepsilon}{2}$, then $(n-1)$ full steps and finally other half-step. A single such step $(n = 1)$ is denoted as follows: 
\begin{align}
    \phi(0) &\rightarrow \phi(\varepsilon) \\
    \pi(0) &\rightarrow \pi(\frac{\varepsilon}{2}) \rightarrow \pi(\varepsilon). 
\end{align}
The simplest integration scheme is second-order accurate: 
\begin{align}\label{eq: leapfrog_scheme}
    \phi(0) &\rightarrow \phi(\varepsilon) = \phi(0) + \pi(\frac{\varepsilon}{2}) \varepsilon \\
    \pi(0) &\rightarrow \pi(\frac{\varepsilon}{2}) = \pi(0) - \frac{\partial S}{\partial \phi}\vert_{\phi(0)}\frac{\varepsilon}{2} \rightarrow \pi(\varepsilon) = \pi(\frac{\varepsilon}{2}) - \frac{\partial S}{\partial \phi} \vert_{\phi(\varepsilon)}\frac{\varepsilon}{2}
\end{align}
When dealing with a series of steps of length $n$, we use the abbreviations $\phi(n \varepsilon) = \phi_n$ and $\pi(n \varepsilon) = \pi_n$. 
\\\\
The invariance of the integration measure can be shown from the product of Jacobians along the sequence: 
\begin{equation}
    \det \left[ \frac{\partial(\pi_1, \phi_1)}{\partial(\pi_0, \phi_0)}\right] = \det \left[ \frac{\partial(\pi_1, \phi_1)}{\partial(\pi_\frac{1}{2}, \phi_1)} \frac{\partial(\pi_\frac{1}{2}, \phi_1)}{\partial(\pi_\frac12, \phi_0)} \frac{\partial(\pi_\frac12, \phi_0)}{\partial(\pi_0, \phi_0)}\right] = 1
\end{equation}
This factorizes into a product of three determinants of upper diagonal matrices with $1$ on the diagonal, hence the result equals $1$ and the integration measure is invariant. 
\\\\
For the reversibility property, use the scheme in Equations \ref{eq: leapfrog_scheme} starting from $(\phi_0, \pi_0)$: 
\begin{align}
    \phi_1 &= \phi_0 + \pi_0 \varepsilon - \frac{1}{2}\frac{\partial S}{\partial \phi} \vert_{\phi_0} \varepsilon^2 \\
    \pi_1 &= \pi_0 - \frac{1}{2} \left( \frac{\partial S}{\partial \phi}\vert_{\phi_0} + \frac{\partial S}{\partial \phi} \vert_{\phi_1}\right) \varepsilon
\end{align}
To go backwards, we start at $(\phi_1, \pi_1)$ and apply the equations of motion \ref{eq: leapfrog_scheme} with a negative step size $-\varepsilon$. This gives
\begin{align}
    \phi(\varepsilon - \varepsilon) &= \phi_1 - \pi_1 \varepsilon - \frac{1}{2}\frac{\partial S}{\partial \phi} \vert_{\phi_1} \varepsilon^2 \\ 
    &= \phi_1 - \pi_1 \varepsilon + \frac{1}{2}\left( \frac{\partial S}{\partial \phi}\vert_{\phi_0} + \frac{\partial S}{\partial \phi} \vert_{\phi_1} \right)\varepsilon^2 - \frac{1}{2}\frac{\partial S}{\partial \phi}\vert_{\phi_1} \varepsilon^2 \\
    &= \phi_1 - \pi_1 \varepsilon + \frac{1}{2} \frac{\partial S}{\partial \phi}\vert_{\phi_0} \varepsilon^2 = \phi_0, 
\end{align}
\begin{align}
    \pi(\varepsilon - \varepsilon) &= \pi_0 = \pi_1 + \left( \frac{\partial S}{\partial \phi} \vert_{\phi_1} + \frac{1}{2}\frac{\partial S}{\partial \phi}\vert_{\phi_0}\right) = \pi_0. 
\end{align}
So the reversibility is proven. It is possible to combine $n$ leapfrog steps. The discretization error for half-steps is $\mathcal{O}(\varepsilon^2)$, for full steps $\mathcal{O}(\varepsilon^3)$ \cite{GattingerLang}. 
\section{The Omelyan2 Integrator}
This integrator is described in great detail in the following papers: \cite{Takaishi2005Symplectic,Omelyan2002Symplectic}. \\\\
The Omelyan2 integrator is a generalization of the standard leapfrog scheme, designed to improve 
energy conservation while preserving the two key properties required for HMC: 
\begin{itemize}
    \item Invariance of the integration measure in the $(\phi, \pi)$ phase space
    \item Reversibility of the trajectory under time reversal
\end{itemize}
Like leapfrog, Omelyan2 integration evolves the fields $\phi$ and conjugate momenta $\pi$ in a sequence of small steps. The key difference is that the momentum and position updates are distributed asymmetrically within each step according to a free parameter $\lambda$, which can be tuned to minimize the discretization error of the Hamiltonian. \\
A single Omelyan step of size $\varepsilon$ can be written as:
\begin{align}
\pi_0 &\rightarrow \pi_{\lambda/2} = \pi_0 - \lambda \varepsilon \frac{\partial S}{\partial \phi}\bigg|_{\phi_0} \\
\phi_0 &\rightarrow \phi_{1/2} = \phi_0 + \frac{\varepsilon}{2} \pi_{\lambda/2} \\
\pi_{\lambda/2} &\rightarrow \pi_{1-\lambda} = \pi_{\lambda/2} - (1-2\lambda) \varepsilon \frac{\partial S}{\partial \phi}\bigg|_{\phi_{1/2}} \\
\phi_{1/2} &\rightarrow \phi_1 = \phi_{1/2} + \frac{\varepsilon}{2} \pi_{1-\lambda} \\
\pi_{1-\lambda} &\rightarrow \pi_1 = \pi_{1-\lambda} - \lambda \varepsilon \frac{\partial S}{\partial \phi}\bigg|_{\phi_1}
\end{align}
Here, $\lambda$ is a free parameter, and the standard leapfrog scheme is recovered for $\lambda = \frac{1}{2}$. In practice, a common 
choice is $\lambda \approx 0.1931833$, which minimizes the error in the Hamiltonian for harmonic systems \cite{Takaishi2005Symplectic,Omelyan2002Symplectic}.\\\\
The Omelyans2 integrator improves on the basic leapfrog scheme in the following ways:
\begin{itemize}
    \item Reduced energy drift: By distributing the momentum updates asymmetrically, the integrator achieves better conservation of the Hamiltonian over long trajectories, which increases the acceptance rate in HMC.  
    \item Preservation of measure and reversibility: Like leapfrog, the Omelyan2 integrator is reversible, ensuring that detailed balance in HMC is maintained.  
    \item Higher efficiency: For a given step size $\varepsilon$, Omelyan2 typically allows longer trajectories or fewer steps while achieving the same or better acceptance, reducing computational cost.  
\end{itemize}
